@ID: ON15D1P1U
@BIDANG: Kombinatorika

Misalkan $n\in\mathbb N$ dan $a,b\in\mathbb Z$ memenuhi

$$
n\mid a^2+b^2+1.
$$

Buktikan bahwa terdapat $x,y\in\mathbb Z$ sehingga

$$
x^2+y^2+1=n.
$$

@ID: ON15D1P2U
@BIDANG: Analisis Real

Misalkan $a,b,c$ bilangan real positif dengan $ab+bc+ca=1$. Buktikan bahwa

$$
\frac{a}{\sqrt{a^2+6}}
+
\frac{b}{\sqrt{b^2+6}}
+
\frac{c}{\sqrt{c^2+6}}
\le1.
$$

@ID: ON15D1P3U
@BIDANG: Analisis Real

Misalkan $f:\mathbb R\to\mathbb R$ fungsi kontinu yang memenuhi

$$
f(x+y)=f(x)f(y),
\qquad \forall x,y\in\mathbb R.
$$

Tentukan semua fungsi $f$ yang memenuhi syarat tersebut.

@ID: ON15D1P4U
@BIDANG: Aljabar Linear

Misalkan $A$ matriks real berukuran $n\times n$ yang memenuhi

$$
A^2=A.
$$

Buktikan bahwa $\operatorname{rank}(A)=\operatorname{tr}(A)$.

@ID: ON15D1P5U
@BIDANG: Kombinatorika

Misalkan terdapat $2n$ orang yang akan dipasangkan menjadi $n$ pasangan. Tentukan banyaknya cara memasangkan mereka.

@ID: ON15D2P1U
@BIDANG: Analisis Real

Misalkan $f:\mathbb R\to\mathbb R$ terdiferensial, $f'(x)>f(x)$ untuk setiap $x\in\mathbb R$, dan $f(x_0)=0$ untuk suatu $x_0\in\mathbb R$.

Buktikan bahwa $f(x)>0$ untuk setiap $x>x_0$.

Sebagai aplikasi, diberikan $c>0$, perlihatkan bahwa persamaan

$$
ce^x=1+x+\frac{x^2}{2}
$$

mempunyai tepat satu akar.

@ID: ON15D2P2U
@BIDANG: Analisis Kompleks

(a) Misalkan $n$ bilangan asli. Buktikan bahwa

$$
\frac{1}{2\pi i}\oint_C\frac{(1+z)^{2n}}{z^{n+1}}\,dz
=
\binom{2n}{n},
$$

dengan $C$ sebarang lingkaran yang mengelilingi titik asal.

(b) Gunakan hasil pada bagian (a) untuk menentukan

$$
\sum_{n=0}^{\infty}\frac{\binom{2n}{n}}{5^n}.
$$

Tunjukkan perhitungannya.

@ID: ON15D2P3U
@BIDANG: Aljabar Linear

Diketahui dua matriks persegi seukuran $M$ dan $N$ yang semua nilai eigennya real positif.

(a) Jika $M=M^\ast$ dan $N=N^\ast$, tunjukkan bahwa semua nilai eigen matriks $MN$ positif.

(b) Berikan contoh matriks $M$ dan $N$ berukuran $3\times3$ sehingga $MN$ hanya mempunyai nilai eigen real negatif. Tunjukkan bahwa contoh yang diberikan memang memenuhi semua persyaratan yang diminta.

@ID: ON15D2P4U
@BIDANG: Struktur Aljabar

(a) Berikan contoh suku banyak $P(x)$ dengan koefisien rasional yang memenuhi

$$
P\!\left(1+2^{2/3}\right)=1+2^{4/3}.
$$

(b) Tentukan semua suku banyak $Q(x)$ dengan koefisien rasional yang memenuhi persamaan pada bagian (a).

@ID: ON15D2P5U
@BIDANG: Kombinatorika

Definisikan

$$
A_n=\frac{1}{n+1}\binom{2n}{n},
\qquad n\in\mathbb N,
$$

dan bentuk fungsi pembangkit

$$
A(x)=\sum_{n\ge0}A_nx^n.
$$

Buktikan bahwa

$$
\sum_{n\ge0}(2n+1)x^nA(x)^{2n+1}
=
\sum_{m\ge0}4^m x^m.
$$